% Secao 5 - O deserto de fluxo (descritivo). Medida primaria F1.

\section{O deserto de fluxo}
\label{sec:fluxo}

O deserto de fluxo substitui ``qual o hospital mais próximo'' por ``a que
distância está o cuidado que o paciente de fato usa''. A medida primária é o
burden efetivo de fluxo F1: a distância de viagem ponderada pelo número de
internações em cada destino, de modo que um município cujos pacientes se espalham
por hubs distantes acumula burden alto mesmo que tenha um hospital ao lado.
Formalmente, para o município $m$,
\[
F1_m \;=\; \frac{\sum_h n_{m,h}\, d\big(m, \text{mun}(h)\big)}{\sum_h n_{m,h}},
\]
onde $n_{m,h}$ é o número de internações de residentes de $m$ no hospital $h$ e
$d(\cdot)$ a distância entre centroides. Operacionalizações alternativas ---
isolamento em uma rede de co-paciência (embeddings) e concentração dos destinos
(HHI) --- concordam com F1 e são tratadas como robustez no apêndice.

A Figura~\ref{fig:flow_map} mostra F1 na mesma escala de quilômetros da
Figura~\ref{fig:distance_map}, e o contraste é o ponto central do relatório. O
burden efetivo mediano é de 49~km --- mais do triplo da distância mediana ao
hospital mais próximo --- e 5\% dos municípios superam 136~km.% src: distribuicao f1_km - p50 49, p75 73, p95 136, max 619
Na mesma paleta, o mapa de fluxo é visivelmente mais quente: vastas áreas do
interior do Nordeste, do Centro-Oeste e do Norte que a distância pintava como
``perto'' carregam burden de viagem alto, porque o cuidado efetivamente usado
está longe. O deserto de fluxo não é um recorte do deserto de distância; é um
mapa próprio, e a Seção~\ref{sec:divergencia} mede exatamente onde os dois
divergem e quem mora na diferença.
